\documentclass[10pt, aspectratio=169]{beamer}
% Preámbulo modular para presentaciones Beamer (Modelación Estadística)
% Diseño y paleta institucional del Tecnológico de Monterrey

\documentclass[aspectratio=169,xcolor={dvipsnames,table}]{beamer}

\usepackage[utf8]{inputenc}
\usepackage[spanish,mexico]{babel}
\usepackage{amsmath,amssymb,amsfonts,mathtools}
\usepackage{graphicx}
\usepackage{listings}
\usepackage{booktabs}
\usepackage{tikz}
\usetikzlibrary{matrix,arrows,decorations.pathmorphing,shapes.geometric,positioning}

% Paleta Institucional Tecnológico de Monterrey
\definecolor{TecAzul}{HTML}{0039A6}       % Azul TEC: color primario institucional
\definecolor{TecAzulOscuro}{HTML}{1A2E51} % Azul marino oscuro
\definecolor{TecRojo}{HTML}{EC2661}       % Rojo de acento (paleta miTec)
\definecolor{TecGris}{HTML}{646464}       % Gris neutro para comentarios y subtítulos
\definecolor{TecFondoBloque}{HTML}{F4F6F9}% Fondo suave para bloques

% Configuración de Tema Beamer
\usetheme{Boadilla}
\usecolortheme[named=TecAzulOscuro]{structure}
\setbeamercolor{alerted text}{fg=TecRojo}
\setbeamercolor{palette primary}{bg=TecAzulOscuro,fg=white}
\setbeamercolor{palette secondary}{bg=TecAzul,fg=white}
\setbeamercolor{palette tertiary}{bg=TecAzulOscuro!80!black,fg=white}
\setbeamercolor{title}{fg=TecAzulOscuro,bg=TecFondoBloque}
\setbeamercolor{frametitle}{fg=TecAzulOscuro,bg=TecFondoBloque}

% Bloques personalizados elegantes
\setbeamercolor{block title}{bg=TecAzulOscuro,fg=white}
\setbeamercolor{block body}{bg=TecFondoBloque,fg=black}
\setbeamercolor{block title alerted}{bg=TecRojo,fg=white}
\setbeamercolor{block body alerted}{bg=TecRojo!8,fg=black}
\setbeamercolor{block title example}{bg=TecAzul,fg=white}
\setbeamercolor{block body example}{bg=TecAzul!8,fg=black}

% Redondear bordes y añadir sombra sutil
\setbeamertemplate{blocks}[rounded][shadow=true]
\setbeamertemplate{navigation symbols}{}

% Pie de página limpio con número de diapositiva
\setbeamertemplate{footline}{
  \leavevmode%
  \hbox{%
  \begin{beamercolorbox}[wd=.7\paperwidth,ht=2.25ex,dp=1ex,left]{author in head/foot}%
    \hspace*{2ex}\insertshortauthor~~(\insertshortinstitute) \hfill \insertshorttitle
  \end{beamercolorbox}%
  \begin{beamercolorbox}[wd=.3\paperwidth,ht=2.25ex,dp=1ex,right]{date in head/foot}%
    \insertshortdate{}\hspace*{2em}
    \insertframenumber{} / \inserttotalframenumber\hspace*{2ex}
  \end{beamercolorbox}}%
  \vskip0pt%
}

% Configuración de Listings de Python para visualización pedagógica de código
\lstset{
    language=Python,
    basicstyle=\ttfamily\footnotesize,
    keywordstyle=\color{TecAzulOscuro}\bfseries,
    stringstyle=\color{TecRojo},
    commentstyle=\color{TecGris}\itshape,
    showstringspaces=false,
    breaklines=true,
    frame=lines,
    rulecolor=\color{TecAzulOscuro!30},
    backgroundcolor=\color{TecFondoBloque},
    aboveskip=0.5em,
    belowskip=0.5em,
    literate={á}{{\'a}}1 {é}{{\'e}}1 {í}{{\'i}}1 {ó}{{\'o}}1 {ú}{{\'u}}1
             {Á}{{\'A}}1 {É}{{\'E}}1 {Í}{{\'I}}1 {Ó}{{\'O}}1 {Ú}{{\'U}}1
             {ñ}{{\~n}}1 {Ñ}{{\~N}}1 {¿}{{?`}}1 {¡}{{!`}}1
}

% Metadatos institucionales por defecto
\title[Modelación Estadística]{Modelación Estadística para Ciencia de Datos e Ingeniería}
\author[J. Castillo Colmenares]{Juliho David Castillo Colmenares}
\institute[Tec de Monterrey]{Tecnológico de Monterrey \\ Escuela de Ingeniería y Ciencias}
\date{\today}

% Comandos y macros estadísticas compartidas para las presentaciones Beamer (Metropolis)

% Conjuntos numéricos básicos
\newcommand{\R}{\mathbb{R}}
\newcommand{\N}{\mathbb{N}}
\newcommand{\Q}{\mathbb{Q}}
\newcommand{\Z}{\mathbb{Z}}
\newcommand{\set}[1]{\left\{ #1 \right\}}
\newcommand{\abs}[1]{\left|#1\right|}
\newcommand{\norm}[1]{\left\| #1 \right\|}

% Comandos estadísticos y probabilísticos del curso
\newcommand{\Var}{\operatorname{Var}}
\newcommand{\cov}{\operatorname{Cov}}
\newcommand{\corr}{\rho}
\newcommand{\comb}[2]{\begin{pmatrix} #1 \\ #2 \end{pmatrix}}
\newcommand{\s}{\sigma}
\newcommand{\particion}{\mathfrak{P}}
\newcommand{\card}[1]{\#\left( #1 \right)}
\newcommand{\seq}[3]{\set{#1}_{#2}^{#3}}
\newcommand{\E}{\mathbb{E}}
\newcommand{\Prob}{\mathbb{P}}

% Entornos pedagógicos y cajas en español académico natural
\newenvironment{discusion}{%
  \begin{alertblock}{Pregunta de Discusión}%
}{%
  \end{alertblock}%
}

\newenvironment{reflexion}{%
  \begin{alertblock}{Reflexión para el Aula}%
}{%
  \end{alertblock}%
}

\newenvironment{paradoja}{%
  \begin{alertblock}{Paradoja de Exploración}%
}{%
  \end{alertblock}%
}

\newenvironment{motivacion}{%
  \begin{exampleblock}{Motivación: Ciencia de Datos en la Práctica}%
}{%
  \end{exampleblock}%
}

\newenvironment{retoclase}{%
  \begin{block}{Reto Rápido en Equipo}%
}{%
  \end{block}%
}


\title{Probability Theory}
\subtitle{Bayes' Theorem and Conditional Inversion}
\author[J. Castillo Colmenares]{Juliho Castillo Colmenares}
\institute{Tecnológico de Monterrey}
\date{}

\begin{document}

% Slide 1: Title Page
\begin{frame}[plain]
  \titlepage
\end{frame}

% Slide 2: Roadmap
\begin{frame}{Chapter 02 Roadmap}
  \begin{itemize}
    \item 02.01 Probability Introduction and Historical Approaches
    \item 02.02 Sample Spaces, Events, and Partitions
    \item 02.03 Kolmogorov's Axioms and Fundamental Theorems
    \item 02.04 Conditional Probability, Multiplication Rule, and Independence
    \item \textbf{\color{TecRojo} 02.05 Bayes' Theorem and Conditional Inversion}
    \item 02.06 Random Sampling (With and Without Replacement)
  \end{itemize}
\end{frame}

% Slide 3: Motivation and Conditional Inversion
\begin{frame}{Motivation: Inverting Conditionals}
  \begin{block}{From Cause to Effect vs. From Effect to Cause}
    In empirical sciences and engineering, we typically model the forward probability of observing an \textbf{effect} $B$ given an underlying \textbf{cause} or hypothesis $A$:
    \begin{align*}
      P(\text{Effect} \mid \text{Cause}) = P(B \mid A) \quad \text{(Likelihood)}
    \end{align*}
    However, in clinical diagnostics, machine learning, and forensic science, the inverse problem is critical: we observe evidence $B$ and must infer the probability of the cause $A$:
    \begin{align*}
      P(\text{Cause} \mid \text{Effect}) = P(A \mid B) \quad \text{(Posterior Probability)}
    \end{align*}
  \end{block}
  \vspace{0.3em}
  \footnotesize
  \emph{Classic example:} We know the probability that a patient with a specific disease tests positive ($P(+ \mid E)$). But clinicians face the inverse problem: given a positive test result, what is the actual probability that the patient has the disease ($P(E \mid +)$)?
\end{frame}

% Slide 4: Deriving Bayes' Theorem (Simple Form)
\begin{frame}{Deriving Bayes' Theorem (Simple Form)}
  If $P(A) > 0$ and $P(B) > 0$, by the definition of conditional probability and the multiplication rule, the joint intersection $A \cap B$ is symmetric:
  \begin{align*}
    P(A \cap B) &= P(A) \cdot P(B \mid A) \\
    P(B \cap A) &= P(B) \cdot P(A \mid B)
  \end{align*}
  Equating both expressions ($P(A \cap B) = P(B \cap A)$):
  \begin{align*}
    P(B) \cdot P(A \mid B) = P(A) \cdot P(B \mid A)
  \end{align*}
  Solving for $P(A \mid B)$ yields the fundamental form of \textbf{Bayes' Theorem}:
  \begin{alertblock}{Bayes' Theorem (Two Hypotheses)}
    \begin{align*}
      P(A \mid B) = \frac{P(A) \cdot P(B \mid A)}{P(B)} = \frac{P(A) \cdot P(B \mid A)}{P(B \mid A)P(A) + P(B \mid A')P(A')}
    \end{align*}
  \end{alertblock}
\end{frame}

% Slide 5: Formal Bayesian Terminology
\begin{frame}{Anatomy and Formal Terminology of Bayes' Theorem}
  In data science and statistical inference, each term in Bayes' formula carries precise epistemological meaning:
  \vspace{0.4em}
  \begin{columns}[T]
    \begin{column}{0.48\textwidth}
      \begin{block}{Prior Probability --- $P(A)$}
        The initial belief, historical baseline, or baseline prevalence of hypothesis $A$ \textbf{before} observing the new evidence $B$.
      \end{block}
      \begin{block}{Likelihood --- $P(B \mid A)$}
        The conditional probability that evidence $B$ occurs under the assumption that hypothesis $A$ is true.
      \end{block}
    \end{column}
    \begin{column}{0.48\textwidth}
      \begin{block}{Marginal Evidence --- $P(B)$}
        The total probability of observing evidence $B$ summed across all possible hypotheses in the system (normalization constant).
      \end{block}
      \begin{block}{Posterior Probability --- $P(A \mid B)$}
        The updated belief regarding hypothesis $A$ \textbf{after} incorporating and processing the observed evidence $B$.
      \end{block}
    \end{column}
  \end{columns}
\end{frame}

% Slide 6: Generalization via Partitions (Chain Rule)
\begin{frame}{Generalized Bayes' Theorem (Exhaustive Partitions)}
  \small
  Let $\{E_1, E_2, \dots, E_N\}$ be an exhaustive ($\bigcup_{i=1}^N E_i = S$) and mutually exclusive ($E_i \cap E_j = \emptyset, i \neq j$) partition of the sample space $S$. For any observable event $B$ with $P(B) > 0$, the Law of Total Probability states:
  \begin{align*}
    P(B) = \sum_{i=1}^N P(B \cap E_i) = \sum_{i=1}^N P(B \mid E_i)P(E_i)
  \end{align*}
  Substituting this denominator into the conditional inversion formula for a specific cell $E_j$:
  \begin{alertblock}{Bayes' Theorem for $N$ States or Hypotheses}
    \begin{align*}
      P(E_j \mid B) = \frac{P(B \mid E_j)P(E_j)}{\sum_{i=1}^N P(B \mid E_i)P(E_i)}, \quad \text{for all } j \in \{1, 2, \dots, N\}
    \end{align*}
  \end{alertblock}
\end{frame}

% Slide 7: Textbook Applied Example (Packaging Machines - Problem Statement)
\begin{frame}{Industrial Example: Packaging Machines (Problem Setup)}
  \footnotesize
  A gelatin manufacturing plant operates three packaging machines with the following production volume shares and defect rates:
  \begin{itemize}
    \item \textbf{Machine 1 ($M_1$):} Accounts for 38\% of production ($P(M_1)=0.38$) with an 11\% defect rate ($P(D \mid M_1)=0.11$).
    \item \textbf{Machine 2 ($M_2$):} Accounts for 32\% of production ($P(M_2)=0.32$) with a 15\% defect rate ($P(D \mid M_2)=0.15$).
    \item \textbf{Machine 3 ($M_3$):} Accounts for 30\% of production ($P(M_3)=0.30$) with a 14\% defect rate ($P(D \mid M_3)=0.14$).
  \end{itemize}
  \vspace{0.4em}
  \begin{block}{Quality Audit Question}
    If management randomly selects one package from the warehouse and it is found to be \textbf{defective ($D$)}, what is the posterior probability that it was packaged by \textbf{Machine 2 ($M_2$)}?
  \end{block}
\end{frame}

% Slide 8: Textbook Applied Example (Solution and Calculation)
\begin{frame}{Industrial Example: Packaging Machines (Solution)}
  \footnotesize
  Applying the generalized Bayes' Theorem by evaluating $M_2$'s contribution relative to total plant defects:
  \begin{align*}
    P(M_2 \mid D) &= \frac{P(D \mid M_2)P(M_2)}{P(D \mid M_1)P(M_1) + P(D \mid M_2)P(M_2) + P(D \mid M_3)P(M_3)} \\[0.6em]
    &= \frac{(0.15)(0.32)}{(0.11)(0.38) + (0.15)(0.32) + (0.14)(0.30)} \\[0.6em]
    &= \frac{0.0480}{0.0418 + 0.0480 + 0.0420} = \frac{0.0480}{0.1318} \approx 0.3642 \text{ (36.42\%)}
  \end{align*}
  \begin{alertblock}{Operational Insight}
    Although Machine 1 processes the largest volume (38\%), Machine 2 is the single primary contributor to overall plant defects ($36.42\%$ of all defective units).
  \end{alertblock}
\end{frame}

% Slide 9: Data Science and Paradoxes (Base Rate Fallacy)
\begin{frame}{Data Science: The Base Rate Fallacy}
  \small
  A widespread cognitive and statistical error involves neglecting the prior probability ($P(A)$) and focusing solely on test sensitivity or precision ($P(B \mid A)$):
  \begin{block}{The False Positive Paradox in Rare Screening}
    If a pathology has a very low prevalence (e.g., $P(E) = 1\%$) and a diagnostic screening test has $95\%$ sensitivity and $90\%$ specificity, the vast majority of positive test results across the general population will be \textbf{false positives}.
  \end{block}
  \vspace{0.3em}
  \footnotesize
  \emph{Mathematical root:} The massive pool of healthy individuals ($99\%$) multiplied by the false positive rate ($10\%$) completely overwhelms the true positives generated from the tiny $1\%$ diseased subpopulation ($0.99 \times 0.10 = 0.099 \gg 0.01 \times 0.95 = 0.0095$).
\end{frame}

% Slide 10: Sequential Bayesian Updating in Machine Learning
\begin{frame}{Sequential Updating: The Engine of Bayesian Learning}
  An elegant computational property of Bayes' Theorem is its dynamic updating mechanism when processing a stream of sequential observations ($B_1, B_2, \dots, B_t$):
  \vspace{0.4em}
  \begin{alertblock}{Today's Posterior Becomes Tomorrow's Prior}
    \begin{align*}
      P(A \mid B_1) &= \frac{P(B_1 \mid A)P(A)}{P(B_1)} \\
      P(A \mid B_1 \cap B_2) &= \frac{P(B_2 \mid A)P(A \mid B_1)}{P(B_2 \mid B_1)}
    \end{align*}
  \end{alertblock}
  \footnotesize
  This recursive structure forms the theoretical backbone of spam filters (\emph{Naïve Bayes}), Kalman tracking filters in robotics, and modern Bayesian reinforcement learning algorithms.
\end{frame}

% Slide 11A: Python Laboratory (Medical Diagnostic Setup)
\begin{frame}[fragile]{Python Laboratory: Medical Test Inversion (1/4)}
  \footnotesize
  False positive tracking and Positive Predictive Value (\texttt{02.05\_bayes\_theorem.py}):
  \vspace{0.2em}
  {\lstset{basicstyle=\ttfamily\fontsize{6.3pt}{7.3pt}\selectfont, aboveskip=0.1em, belowskip=0.1em}
  \lstinputlisting[firstline=18, lastline=38]{../../code/02_teoria_probabilidad/02.05_bayes_theorem.py}}
\end{frame}

% Slide 11B: Python Laboratory (Naive Bayes Anti-Spam Filter)
\begin{frame}[fragile]{Python Laboratory: Naïve Bayes Spam Filter (2/4)}
  \footnotesize
  Exact posterior evaluation of a text classifier under conditional independence:
  \vspace{0.2em}
  {\lstset{basicstyle=\ttfamily\fontsize{6.3pt}{7.3pt}\selectfont, aboveskip=0.1em, belowskip=0.1em}
  \lstinputlisting[firstline=43, lastline=57]{../../code/02_teoria_probabilidad/02.05_bayes_theorem.py}}
\end{frame}

% Slide 11C: Python Laboratory (Sequential Updating)
\begin{frame}[fragile]{Python Laboratory: Sequential Updating (3/4)}
  \footnotesize
  Iterative evolution of prior $\to$ posterior beliefs across successive evidence:
  \vspace{0.2em}
  {\lstset{basicstyle=\ttfamily\fontsize{6.3pt}{7.3pt}\selectfont, aboveskip=0.1em, belowskip=0.1em}
  \lstinputlisting[firstline=66, lastline=86]{../../code/02_teoria_probabilidad/02.05_bayes_theorem.py}}
\end{frame}

% Slide 11D: Python Laboratory (Console Output)
\begin{frame}[fragile]{Python Laboratory: Terminal Output (4/4)}
  \footnotesize
  Empirical convergence and verification across all three computational experiments:
  \vspace{0.1em}
  \begin{block}{Standard Terminal Output}
    \fontsize{5.7pt}{6.6pt}\selectfont
    \begin{verbatim}
--- 1. Medical Diagnostic Test (Base Rate Fallacy) ---
Total Test Positives:  11,667 / 100,000
True Positives:        1,891 (True Diseased among Positives)
Empirical PPV (Sim):   0.1621 (16.21%) | Theoretical PPV: 0.1624 (16.24%)

--- 2. Naive Bayes Spam Filter (Keywords W1, W2, W3) ---
Joint Likelihood P(W|Spam):  0.33600 | Joint Likelihood P(W|Legit): 0.00200
Marginal Evidence P(W):      0.20240 | Posterior P(Spam | W): 0.996047 (99.60%)

--- 3. Sequential Bayesian Updating (Coin Flips: H, H, T, H, H) ---
Step 0 | P(Biased Coin = 0.80): 0.5000 (50.00%)
Step 1 | P(Biased Coin = 0.80): 0.6154 (61.54%)
Step 2 | P(Biased Coin = 0.80): 0.7191 (71.91%)
Step 3 | P(Biased Coin = 0.80): 0.5059 (50.59%)
Step 4 | P(Biased Coin = 0.80): 0.6210 (62.10%)
Step 5 | P(Biased Coin = 0.80): 0.7239 (72.39%)
    \end{verbatim}
  \end{block}
\end{frame}

% Slide 12: Textbook Exercises — Tier 1: Fundamental
\begin{frame}{Textbook Exercises — Tier 1: Fundamental (Direct)}
  \footnotesize
  \textbf{Problem 2.5.1 (Die Roll and Parity Conditional):}
  A balanced 6-sided die is rolled once. Let $A=\{1,2,3\}$ be the event of rolling a number strictly less than 4, and $I=\{1,3,5\}$ be the event of rolling an odd number.
  \begin{itemize}
    \item Without prior information: $P(A) = \frac{3}{6} = 0.50$.
    \item Given that the roll resulted in an odd number ($I$):
      \begin{align*}
        P(A \mid I) = \frac{P(A \cap I)}{P(I)} = \frac{\card{\{1,3\}}}{\card{\{1,3,5\}}} = \frac{2}{3} \approx 0.6667 \text{ (66.67\%)}
      \end{align*}
  \end{itemize}
  \vspace{0.2em}
  \textbf{Problem 2.5.6 (Left-Handed Students by Gender):}
  In a class, $60\%$ are women ($M$) and $40\%$ are men ($H$). Left-handed rates are $P(Z \mid M)=0.10$ and $P(Z \mid H)=0.15$. If a randomly selected student is left-handed ($Z$), the posterior probability of being female is:
  \begin{align*}
    P(M \mid Z) = \frac{P(Z \mid M)P(M)}{P(Z \mid M)P(M) + P(Z \mid H)P(H)} = \frac{(0.10)(0.60)}{(0.10)(0.60) + (0.15)(0.40)} = \frac{0.06}{0.12} = 50\%
  \end{align*}
\end{frame}

% Slide 13: Textbook Exercises — Tier 1 (Cont.) & Tier 2: Operational
\begin{frame}{Textbook Exercises — Tier 1 (Cont.) and Tier 2: Operational}
  \footnotesize
  \textbf{Problem 2.5.4 (Study Habits and Exam Passage):}
  A student studies diligently with probability $P(E)=0.70$. The probability of passing ($A$) given studying is $P(A \mid E)=0.90$, while without studying it is $P(A \mid E')=0.30$.
  \begin{itemize}
    \item Total probability of passing: $P(A) = (0.90)(0.70) + (0.30)(0.30) = 0.63 + 0.09 = 0.72$ ($72\%$).
    \item If the student passed, posterior probability of studying: $P(E \mid A) = \frac{0.63}{0.72} = \frac{7}{8} = 87.5\%$.
    \item If the student failed ($A'$), probability of not studying: $P(E' \mid A') = \frac{(0.70)(0.30)}{1 - 0.72} = \frac{0.21}{0.28} = 75\%$.
  \end{itemize}
  \vspace{0.2em}
  \textbf{Problem 2.5.7 (Urns and Two-Dice Sums):}
  Box 1: 3 red, 2 blue ($P(R \mid C_1)=\frac{3}{5}$). Box 2: 2 red, 8 blue ($P(R \mid C_2)=\frac{1}{5}$). Box 1 is chosen if the sum of two dice is $\le 6$ ($P(C_1)=\frac{15}{36}=\frac{5}{12}$, $P(C_2)=\frac{7}{12}$).
  If a red marble ($R$) was drawn, the posterior probability of Box 1 is:
  \begin{align*}
    P(C_1 \mid R) = \frac{(3/5)(5/12)}{(3/5)(5/12) + (1/5)(7/12)} = \frac{15/60}{22/60} = \frac{15}{22} \approx 68.18\%
  \end{align*}
\end{frame}

% Slide 14: Textbook Exercises — Tier 2 (Cont.)
\begin{frame}{Textbook Exercises — Tier 2: Operational (Intermediate)}
  \footnotesize
  \textbf{Problem 2.5.2 (Clinical Diagnostic Screening and PPV/NPV):}
  A pathology has prevalence $P(E)=0.02$. The screening test has Sensitivity $P(+ \mid E)=0.95$ and Specificity $P(- \mid E')=0.90$ (hence false positive rate $P(+ \mid E')=0.10$).
  \begin{itemize}
    \item Total probability of testing positive: $P(+) = (0.95)(0.02) + (0.10)(0.98) = 0.019 + 0.098 = 0.117$ ($11.7\%$).
    \item \textbf{Positive Predictive Value (PPV):} $P(E \mid +) = \frac{0.019}{0.117} \approx 16.24\%$. Notice that $83.76\%$ of positives are false alarms!
    \item \textbf{Negative Predictive Value (NPV):} $P(E' \mid -) = \frac{(0.90)(0.98)}{1 - 0.117} = \frac{0.882}{0.883} \approx 99.89\%$.
  \end{itemize}
  \vspace{0.2em}
  \textbf{Problem 2.5.5 (Noisy Binary Digital Channel):}
  A channel transmits $60\%$ 0-bits ($T_0$) and $40\%$ 1-bits ($T_1$). Due to noise, $P(R_1 \mid T_0)=0.02$ and $P(R_1 \mid T_1)=0.97$.
  \begin{itemize}
    \item Probability of receiving a 1-bit: $P(R_1) = (0.02)(0.60) + (0.97)(0.40) = 0.012 + 0.388 = 0.40$ ($40\%$).
    \item If a 1-bit is received, posterior transmission certainty: $P(T_1 \mid R_1) = \frac{0.388}{0.400} = 97.00\%$.
  \end{itemize}
\end{frame}

% Slide 15: Textbook Exercises — Tier 3: Analytical
\begin{frame}{Textbook Exercises — Tier 3: Analytical (Advanced)}
  \footnotesize
  \textbf{Problem 2.5.3 (Semiconductor Assembly Lines):}
  Three lines contribute $P(M_1)=0.50$, $P(M_2)=0.30$, $P(M_3)=0.20$ of production. Defect rates are $P(D \mid M_1)=0.03$, $P(D \mid M_2)=0.05$, $P(D \mid M_3)=0.07$. Marginal defect probability:
  $P(D) = (0.03)(0.50) + (0.05)(0.30) + (0.07)(0.20) = 0.015 + 0.015 + 0.014 = 0.044$.
  \begin{itemize}
    \item Posterior vector: $P(M_1 \mid D) = \frac{15}{44} \approx 34.09\%$, $P(M_2 \mid D) = \frac{15}{44} \approx 34.09\%$, $P(M_3 \mid D) = \frac{14}{44} \approx 31.82\%$.
    \item \emph{Insight:} Both $M_1$ and $M_2$ are tied as the most probable causes of a defect, because $M_1$'s massive production share exacts the same absolute defect count as $M_2$'s higher defect rate.
  \end{itemize}
  \vspace{0.2em}
  \textbf{Problem 2.5.8 (Rigorous Proof for Finite Partitions):}
  Given an exhaustive and mutually exclusive partition $\{B_1, \dots, B_k\}$, for any state $j$:
  $P(B_j \cap A) = P(A \mid B_j)P(B_j)$. Since $A = \bigcup_{i=1}^k (A \cap B_i)$ is a disjoint union, by Kolmogorov's 3rd Axiom, $P(A) = \sum_{i=1}^k P(A \mid B_i)P(B_i)$. Substituting into $P(B_j \mid A) = \frac{P(B_j \cap A)}{P(A)}$ yields Bayes' Theorem.
\end{frame}

% Slide 16: Textbook Exercises — Tier 4: Challenging
\begin{frame}{Textbook Exercises — Tier 4: Challenging (Expert)}
  \footnotesize
  \textbf{Problem 2.5.9 (Prosecutor's Fallacy and DNA Odds Ratio):}
  Crime scene DNA matches a suspect ($E$). Random population match rate ($N=5,000,000$) is $\alpha=10^{-6}$. The prosecutor asserts $99.9999\%$ guilt certainty.
  \begin{itemize}
    \item Prior guilt odds ($I$): $\frac{P(I)}{P(I')} = \frac{1 / 5,000,000}{1} = 2 \times 10^{-7}$.
    \item Likelihood ratio: $\text{LR} = \frac{P(E \mid I)}{P(E \mid I')} = \frac{1}{10^{-6}} = 1,000,000$.
    \item Posterior odds: $\text{Odds}(I \mid E) = 1,000,000 \times (2 \times 10^{-7}) = 0.20 = \frac{1}{5} \implies P(I \mid E) = \frac{0.20}{1.20} = \frac{1}{6} \approx 16.67\%$.
  \end{itemize}
  \vspace{0.1em}
  \textbf{Problem 2.5.10 (Naïve Bayes Spam Classifier Decision Rule):}
  Server processes $60\%$ spam ($S$) and $40\%$ legitimate ($S'$). Keywords $W = W_1 \cap W_2 \cap W_3$ have joint conditional likelihoods: $P(W \mid S) = (0.80)(0.70)(0.60) = 0.336$, and $P(W \mid S') = (0.10)(0.20)(0.10) = 0.002$.
  \begin{align*}
    P(S \mid W) = \frac{(0.336)(0.60)}{(0.336)(0.60) + (0.002)(0.40)} = \frac{0.2016}{0.2016 + 0.0008} = \frac{0.2016}{0.2024} \approx 99.60\%
  \end{align*}
\end{frame}

% Slide 17: Synthesis and Next Session Challenge
\begin{frame}{Synthesis and Challenge for the Next Session}
  \begin{block}{Core Conclusions from Bayes' Theorem}
    \begin{itemize}
      \item \textbf{Rigorous Inversion:} Bayes provides the mathematical bridge linking observed empirical effects to their underlying generative hypotheses.
      \item \textbf{The Weight of the Prior:} Disregarding baseline marginal prevalence ($P(A)$) leads directly to severe diagnostic fallacies, especially in rare screening scenarios.
      \item \textbf{Dynamic Learning:} Sequential Bayesian updating allows intelligent systems to continuously refine their beliefs under accumulating uncertainty.
    \end{itemize}
  \end{block}
  \vspace{0.4em}
  \begin{alertblock}{Challenge for Session 02.06: Random Sampling}
    What happens when we analyze sequences of observations drawn without replacement from finite populations? We will study the structural transition between exact hypergeometric dependence and binomial approximations.
  \end{alertblock}
\end{frame}

\end{document}
